\chapter{第5章：维数公式 习题}
\difficultykey

\section{基础题 \easy}

\begin{problem}{低权维数表}
\easy
对偶数 $k=0,2,4,\dots,24$，计算 $M_k(\SL_2(\Z))$ 的维数。
\end{problem}
\begin{solution}
精确维数公式是
\[
\dim M_k(\SL_2(\Z))=
\begin{cases}
1,&k=0,\\
0,&k=2,\\
\left\lfloor \dfrac{k}{12}\right\rfloor+1,&k\ge 4\text{ 为偶数且 }k\not\equiv 2\pmod{12},\\
\left\lfloor \dfrac{k}{12}\right\rfloor,&k\ge 4\text{ 为偶数且 }k\equiv 2\pmod{12}.
\end{cases}
\]
逐项代入得到
\[
\begin{array}{c|ccccccccccccc}
k&0&2&4&6&8&10&12&14&16&18&20&22&24\\ \hline
\dim M_k&1&0&1&1&1&1&2&1&2&2&2&2&3
\end{array}
\]
这就是所求表格。
\end{solution}

\begin{problem}{低权尖点维数}
\easy
对偶数 $k=0,2,4,\dots,24$，计算 $S_k(\SL_2(\Z))$ 的维数。
\end{problem}
\begin{solution}
对 $k\ge 4$ 的偶数，有
\[
\dim S_k(\SL_2(\Z))=\dim M_k(\SL_2(\Z))-1,
\]
而 $S_0=S_2=0$。由上一题立刻得到
\[
\begin{array}{c|ccccccccccccc}
k&0&2&4&6&8&10&12&14&16&18&20&22&24\\ \hline
\dim S_k&0&0&0&0&0&0&1&0&1&1&1&1&2
\end{array}
\]
特别地，权 $12$ 才第一次出现非零尖点形式。
\end{solution}

\begin{problem}{价公式检验}
\easy
用价公式分别验证 $E_4$、$E_6$、$\Delta$ 的零点分布：
$E_4$ 在 $\rho=e^{2\pi i/3}$ 处有一个简单零点，$E_6$ 在 $i$ 处有一个简单零点，$\Delta$ 在尖点 $\infty$ 处有一个简单零点。
\end{problem}
\begin{solution}
对非零 $f\in M_k(\SL_2(\Z))$，价公式是
\[
\ord_\infty(f)+\frac{1}{2}\ord_i(f)+\frac{1}{3}\ord_\rho(f)+
\sum_{p\ne i,\rho,\infty}\ord_p(f)=\frac{k}{12}.
\]

对 $E_4$，右边等于 $4/12=1/3$。因为 $E_4$ 不是尖点形式，故 $\ord_\infty(E_4)=0$；又各项都非负，所以只能是
\[
\frac{1}{3}\ord_\rho(E_4)=\frac{1}{3},
\]
即 $\ord_\rho(E_4)=1$，其余点无零点。

对 $E_6$，右边等于 $6/12=1/2$。同理 $\ord_\infty(E_6)=0$，于是只能有
\[
\frac{1}{2}\ord_i(E_6)=\frac{1}{2},
\]
即 $\ord_i(E_6)=1$，其余点无零点。

对 $\Delta$，右边等于 $12/12=1$。而 $\Delta(q)=q+O(q^2)$，所以 $\ord_\infty(\Delta)=1$。这样左边已经达到 $1$，故上半平面内没有别的零点。
\end{solution}

\begin{problem}{低权空间的显式基}
\easy
利用 $M_*(\SL_2(\Z))=\C[E_4,E_6]$，分别写出 $M_8$、$M_{10}$、$M_{12}$、$M_{14}$ 的一组基。
\end{problem}
\begin{solution}
只要解方程 $4a+6b=k$，再取单项式 $E_4^aE_6^b$ 即可。

当 $k=8$ 时，唯一解是 $(a,b)=(2,0)$，故
\[
M_8=\C E_4^2.
\]
当 $k=10$ 时，唯一解是 $(1,1)$，故
\[
M_{10}=\C E_4E_6.
\]
当 $k=12$ 时，解为 $(3,0)$ 与 $(0,2)$，故
\[
M_{12}=\C E_4^3\oplus \C E_6^2.
\]
当 $k=14$ 时，唯一解是 $(2,1)$，故
\[
M_{14}=\C E_4^2E_6.
\]
这四个空间的维数分别是 $1,1,2,1$，与维数公式完全一致。
\end{solution}

\begin{problem}{权八恒等式}
\easy
证明 $E_4^2=E_8$。
\end{problem}
\begin{solution}
$E_4^2$ 与 $E_8$ 都是权 $8$ 的模形式，而且它们的常数项都等于 $1$。
由维数公式，$M_8(\SL_2(\Z))$ 是一维空间，所以任何两个常数项同为 $1$ 的权 $8$ 模形式必然相等。
因此
\[
E_4^2=E_8.
\]
\end{solution}

\begin{problem}{小水平的权二空间}
\easy
计算 $\dim S_2(\Gamma_0(N))$，其中 $N=11,17,19$。
\end{problem}
\begin{solution}
第 4 章已经给出
\[
\dim S_2(\Gamma_0(N))=g\bigl(X_0(N)\bigr).
\]
而小水平表中
\[
g\bigl(X_0(11)\bigr)=g\bigl(X_0(17)\bigr)=g\bigl(X_0(19)\bigr)=1.
\]
所以
\[
\dim S_2(\Gamma_0(11))=
\dim S_2(\Gamma_0(17))=
\dim S_2(\Gamma_0(19))=1.
\]
这说明这三个级别上都恰好出现一条权 $2$ 新形式方向。
\end{solution}

\begin{problem}{更多权二例子}
\easy
计算 $\dim S_2(\Gamma_0(N))$，其中 $N=14,15,20$。
\end{problem}
\begin{solution}
同样利用
\[
\dim S_2(\Gamma_0(N))=g\bigl(X_0(N)\bigr).
\]
小水平表给出
\[
g\bigl(X_0(14)\bigr)=g\bigl(X_0(15)\bigr)=g\bigl(X_0(20)\bigr)=1.
\]
因此
\[
\dim S_2(\Gamma_0(14))=
\dim S_2(\Gamma_0(15))=
\dim S_2(\Gamma_0(20))=1.
\]
所以这三个级别也都是“一维权二尖点空间”的例子。
\end{solution}

\begin{problem}{尖点较多时的权二维数}
\easy
已知对 $\Gamma_0(12)$ 有
\[
g=0,\qquad \nu_2=0,\qquad \nu_3=0,\qquad \nu_\infty=6.
\]
计算 $\dim M_2(\Gamma_0(12))$ 与 $\dim S_2(\Gamma_0(12))$。
\end{problem}
\begin{solution}
一般维数公式在 $k=2$ 时给出
\[
\dim M_2(\Gamma)= (2-1)(g-1)+\left\lfloor\frac{2}{4}\right\rfloor\nu_2+
\left\lfloor\frac{2}{3}\right\rfloor\nu_3+\frac{2}{2}\nu_\infty.
\]
代入 $g=0$、$\nu_2=\nu_3=0$、$\nu_\infty=6$，得到
\[
\dim M_2(\Gamma_0(12))=-1+0+0+6=5.
\]
另一方面，
\[
\dim S_2(\Gamma_0(12))=g=0.
\]
所以答案是
\[
\dim M_2(\Gamma_0(12))=5,
\qquad
\dim S_2(\Gamma_0(12))=0.
\]
这里全部五维都来自 Eisenstein 部分。
\end{solution}

\begin{problem}{全模群上不存在权二形式}
\easy
证明 $M_2(\SL_2(\Z))=0$。
\end{problem}
\begin{solution}
可以直接用一般维数公式。对 $\SL_2(\Z)$，有
\[
g=0,\qquad \nu_2=1,\qquad \nu_3=1,\qquad \nu_\infty=1.
\]
于是
\[
\dim M_2(\SL_2(\Z))=(2-1)(-1)+\left\lfloor\frac{2}{4}\right\rfloor+
\left\lfloor\frac{2}{3}\right\rfloor+1=-1+0+0+1=0.
\]
所以 $M_2(\SL_2(\Z))$ 只能是零空间。
\end{solution}

\section{进阶题 \medium}

\begin{problem}{单项式计数推出维数公式}
\medium
只用分次环结构 $M_*(\SL_2(\Z))=\C[E_4,E_6]$，重新推出 $M_k(\SL_2(\Z))$ 的精确维数公式。
\end{problem}
\begin{solution}
权 $k$ 的单项式基候选正是满足
\[
4a+6b=k,
\qquad a,b\ge 0,
\]
的所有 $E_4^aE_6^b$。把等式除以 $2$，记 $n=k/2$，就要数
\[
2a+3b=n
\]
的非负整数解个数。

设 $n=6m+r$，其中 $r\in\{0,1,2,3,4,5\}$。由 $2a=n-3b$ 可知 $b$ 与 $n$ 同奇偶。又 $0\le b\le \lfloor n/3\rfloor$。
于是可取的 $b$ 的个数恰好是：
\[
\#\{b\}=\begin{cases}
m+1,&r\neq 1,\\
m,&r=1.
\end{cases}
\]
而 $r=1$ 正好等价于 $k=2n\equiv 2\pmod{12}$。
因此
\[
\dim M_k(\SL_2(\Z))=
\begin{cases}
\left\lfloor\dfrac{k}{12}\right\rfloor+1,&k\not\equiv 2\pmod{12},\\
\left\lfloor\dfrac{k}{12}\right\rfloor,&k\equiv 2\pmod{12}.
\end{cases}
\]
这就是精确维数公式。
\end{solution}

\begin{problem}{尖点形式都被判别式整除}
\medium
证明对偶数 $k\ge 12$，乘法映射
\[
M_{k-12}(\SL_2(\Z))\longrightarrow S_k(\SL_2(\Z)),
\qquad g\longmapsto \Delta g
\]
是同构。
\end{problem}
\begin{solution}
这个映射显然线性，而且因为 $\Delta$ 是权 $12$ 尖点形式，所以 $\Delta g$ 一定落在 $S_k$ 中。

单射显然：若 $\Delta g=0$，由于 $\Delta\neq 0$，只能有 $g=0$。

为证满射，取任意 $f\in S_k(\SL_2(\Z))$。由于 $f$ 是尖点形式，它在 $\infty$ 处至少有一阶零点；而 $\Delta$ 在 $\infty$ 处恰有一阶零点，并且在上半平面内无零点，所以
\[
\frac{f}{\Delta}
\]
在整个模曲线上全纯，且权为 $k-12$。因此 $f/\Delta\in M_{k-12}(\SL_2(\Z))$，即 $f=\Delta g$。

故该映射为同构，从而
\[
S_k(\SL_2(\Z))=\Delta\, M_{k-12}(\SL_2(\Z)).
\]
\end{solution}

\begin{problem}{若干尖点空间的基}
\medium
求 $S_{16}(\SL_2(\Z))$、$S_{18}(\SL_2(\Z))$、$S_{24}(\SL_2(\Z))$ 的一组基。
\end{problem}
\begin{solution}
由上一题，
\[
S_k(\SL_2(\Z))=\Delta\,M_{k-12}(\SL_2(\Z)).
\]

对 $k=16$，有 $M_4=\C E_4$，所以
\[
S_{16}=\C\,\Delta E_4.
\]

对 $k=18$，有 $M_6=\C E_6$，所以
\[
S_{18}=\C\,\Delta E_6.
\]

对 $k=24$，有 $M_{12}=\C E_4^3\oplus \C E_6^2$，所以
\[
S_{24}=\C\,\Delta E_4^3\oplus \C\,\Delta E_6^2.
\]
利用 $E_4^3-E_6^2=1728\Delta$，还可改写成另一组基
\[
S_{24}=\C\,\Delta E_4^3\oplus \C\,\Delta^2.
\]
\end{solution}

\begin{problem}{判别式恒等式}
\medium
证明存在常数 $c\in\C$ 使
\[
E_4^3-E_6^2=c\Delta,
\]
并求出 $c$。
\end{problem}
\begin{solution}
$E_4^3-E_6^2$ 是权 $12$ 模形式。它的常数项是 $1-1=0$，所以它在尖点 $\infty$ 处消失，因而属于 $S_{12}(\SL_2(\Z))$。

但 $S_{12}(\SL_2(\Z))$ 是一维空间，因此必有
\[
E_4^3-E_6^2=c\Delta
\]
对某个常数 $c$ 成立。

比较 $q$-展开首项即可求 $c$。由
\[
E_4=1+240q+O(q^2),
\qquad
E_6=1-504q+O(q^2),
\qquad
\Delta=q+O(q^2),
\]
得到
\[
E_4^3-E_6^2=(1+720q+O(q^2))-(1-1008q+O(q^2))=1728q+O(q^2).
\]
所以 $c=1728$，即
\[
E_4^3-E_6^2=1728\Delta.
\]
\end{solution}

\begin{problem}{两个素数水平的亏格}
\medium
利用
\[
g\bigl(X_0(p)\bigr)=1+\frac{p+1}{12}-\frac{\nu_2}{4}-\frac{\nu_3}{3}-1
\]
计算 $g(X_0(23))$ 与 $g(X_0(37))$。
\end{problem}
\begin{solution}
对素数 $p>3$，有 $\mu=p+1$，且尖点个数 $\nu_\infty=2$，因此公式可写成
\[
g\bigl(X_0(p)\bigr)=\frac{p+1}{12}-\frac{\nu_2}{4}-\frac{\nu_3}{3}.
\]

对 $p=23$，因为 $23\equiv 3\pmod 4$ 且 $23\equiv 2\pmod 3$，所以
\[
\nu_2=0,\qquad \nu_3=0.
\]
于是
\[
g(X_0(23))=\frac{24}{12}=2.
\]

对 $p=37$，因为 $37\equiv 1\pmod 4$ 且 $37\equiv 1\pmod 3$，所以
\[
\nu_2=2,\qquad \nu_3=2.
\]
于是
\[
g(X_0(37))=\frac{38}{12}-\frac{2}{4}-\frac{2}{3}=\frac{19}{6}-\frac{1}{2}-\frac{2}{3}=2.
\]
因此
\[
g(X_0(23))=g(X_0(37))=2.
\]
\end{solution}

\begin{problem}{一级别十一的高权维数}
\medium
对 $\Gamma_0(11)$，已知
\[
g=1,\qquad \nu_2=0,\qquad \nu_3=0,\qquad \nu_\infty=2.
\]
计算 $\dim M_4(\Gamma_0(11))$、$\dim S_4(\Gamma_0(11))$、$\dim M_6(\Gamma_0(11))$、$\dim S_6(\Gamma_0(11))$。
\end{problem}
\begin{solution}
一般维数公式给出
\[
\dim M_k(\Gamma)=(k-1)(g-1)+\left\lfloor\frac{k}{4}\right\rfloor\nu_2+
\left\lfloor\frac{k}{3}\right\rfloor\nu_3+\frac{k}{2}\nu_\infty,
\]
以及对偶数 $k\ge 4$，
\[
\dim S_k(\Gamma)=(k-1)(g-1)+\left\lfloor\frac{k}{4}\right\rfloor\nu_2+
\left\lfloor\frac{k}{3}\right\rfloor\nu_3+\left(\frac{k}{2}-1\right)\nu_\infty.
\]
代入 $g=1$、$\nu_2=\nu_3=0$、$\nu_\infty=2$，得到
\[
\dim M_k(\Gamma_0(11))=k,
\qquad
\dim S_k(\Gamma_0(11))=k-2\quad (k\ge 4).
\]
于是
\[
\dim M_4=4,
\quad \dim S_4=2,
\quad \dim M_6=6,
\quad \dim S_6=4.
\]
\end{solution}

\begin{problem}{椭圆点个数与权二维数}
\medium
对素数 $p=13,17,19,23$，计算 $\nu_2$、$\nu_3$ 与 $\dim S_2(\Gamma_0(p))$。
\end{problem}
\begin{solution}
对素数 $p>3$，有
\[
\nu_2=1+\left(\frac{-1}{p}\right),
\qquad
\nu_3=1+\left(\frac{-3}{p}\right),
\qquad
\dim S_2(\Gamma_0(p))=g(X_0(p)).
\]

逐个计算：
\begin{itemize}
  \item $p=13$：$13\equiv 1\pmod4$，$13\equiv 1\pmod3$，故 $\nu_2=2$、$\nu_3=2$，从而
  \[
  g=\frac{14}{12}-\frac{2}{4}-\frac{2}{3}=0.
  \]
  所以 $\dim S_2(\Gamma_0(13))=0$。

  \item $p=17$：$17\equiv 1\pmod4$，$17\equiv 2\pmod3$，故 $\nu_2=2$、$\nu_3=0$，从而
  \[
  g=\frac{18}{12}-\frac{2}{4}=1.
  \]
  所以 $\dim S_2(\Gamma_0(17))=1$。

  \item $p=19$：$19\equiv 3\pmod4$，$19\equiv 1\pmod3$，故 $\nu_2=0$、$\nu_3=2$，从而
  \[
  g=\frac{20}{12}-\frac{2}{3}=1.
  \]
  所以 $\dim S_2(\Gamma_0(19))=1$。

  \item $p=23$：$23\equiv 3\pmod4$，$23\equiv 2\pmod3$，故 $\nu_2=0$、$\nu_3=0$，从而
  \[
  g=\frac{24}{12}=2.
  \]
  所以 $\dim S_2(\Gamma_0(23))=2$。
\end{itemize}
\end{solution}

\begin{problem}{一维尖点空间中的 Hecke 基}
\medium
设 $\dim S_k(\Gamma)=1$。证明该空间中的唯一归一化向量自动是所有 Hecke 算子的共同特征向量。
\end{problem}
\begin{solution}
设 $S_k(\Gamma)=\C f$，其中 $f$ 的首项归一化为 $q+O(q^2)$。对任意 Hecke 算子 $T_n$，因为 $T_n$ 保持 $S_k(\Gamma)$，所以 $T_n f$ 仍落在这一维空间里。故存在常数 $\lambda_n$ 使
\[
T_n f=\lambda_n f.
\]
这正说明 $f$ 是 $T_n$ 的特征向量。

由于对每个 $n$ 都成立，所以 $f$ 是全部 Hecke 算子的共同特征向量。也就是说，一维尖点空间自动已经是 Hecke 基。
\end{solution}

\begin{problem}{尖点处的最大消失阶}
\medium
设 $0\ne f\in M_k(\SL_2(\Z))$，其中 $k$ 为偶数。证明：
\[
\ord_\infty(f)\le \left\lfloor\frac{k}{12}\right\rfloor,
\]
并说明当 $k\equiv 2\pmod{12}$ 时实际上有更强的不等式
\[
\ord_\infty(f)\le \left\lfloor\frac{k}{12}\right\rfloor-1.
\]
\end{problem}
\begin{solution}
先证第一条。价公式左边每一项都非负，所以
\[
\ord_\infty(f)\le \frac{k}{12}.
\]
由于左边是整数，便得
\[
\ord_\infty(f)\le \left\lfloor\frac{k}{12}\right\rfloor.
\]

再看 $k\equiv 2\pmod{12}$ 的情形。写 $k=12m+2$。若存在非零 $f$ 满足 $\ord_\infty(f)\ge m$，则 $f/\Delta^m$ 是权 $2$ 的全纯模形式；但 $M_2(\SL_2(\Z))=0$，矛盾。故只能有
\[
\ord_\infty(f)\le m-1=\left\lfloor\frac{k}{12}\right\rfloor-1.
\]
这说明余类 $2\pmod{12}$ 恰好比其他余类少掉一个可用的 $\Delta$ 因子。
\end{solution}

\section{挑战题 \hard}

\begin{problem}{从一般公式恢复全模群公式}
\hard
从一般维数公式
\[
\dim M_k(\Gamma)=(k-1)(g-1)+\left\lfloor\frac{k}{4}\right\rfloor\nu_2+
\left\lfloor\frac{k}{3}\right\rfloor\nu_3+\frac{k}{2}\nu_\infty
\]
出发，推出 $\SL_2(\Z)$ 的精确维数公式。
\end{problem}
\begin{solution}
对 $\Gamma=\SL_2(\Z)$，有
\[
g=0,
\qquad
\nu_2=1,
\qquad
\nu_3=1,
\qquad
\nu_\infty=1.
\]
因此
\[
\dim M_k(\SL_2(\Z))=-(k-1)+\left\lfloor\frac{k}{4}\right\rfloor+
\left\lfloor\frac{k}{3}\right\rfloor+\frac{k}{2}.
\]
现在把偶数 $k$ 按模 $12$ 分六类讨论。写 $k=12m+r$，其中
$r\in\{0,2,4,6,8,10\}$。

逐类计算可得：
\begin{itemize}
  \item 若 $r=0,4,6,8,10$，则
  \[
  -(k-1)+\left\lfloor\frac{k}{4}\right\rfloor+
  \left\lfloor\frac{k}{3}\right\rfloor+\frac{k}{2}=m+1.
  \]
  \item 若 $r=2$，则
  \[
  -(k-1)+\left\lfloor\frac{k}{4}\right\rfloor+
  \left\lfloor\frac{k}{3}\right\rfloor+\frac{k}{2}=m.
  \]
\end{itemize}
于是
\[
\dim M_k(\SL_2(\Z))=
\begin{cases}
\left\lfloor\dfrac{k}{12}\right\rfloor+1,&k\not\equiv 2\pmod{12},\\
\left\lfloor\dfrac{k}{12}\right\rfloor,&k\equiv 2\pmod{12}.
\end{cases}
\]
这正是全模群的精确维数公式。
\end{solution}

\begin{problem}{尖点理想由判别式生成}
\hard
证明分次理想
\[
S_*(\SL_2(\Z))=\bigoplus_{k\ge 0} S_k(\SL_2(\Z))
\]
由单个元素 $\Delta$ 生成。
\end{problem}
\begin{solution}
要证的是：每个尖点形式都唯一写成 $\Delta g$ 的形式，其中 $g$ 是某个较低权模形式。

设 $f\in S_k(\SL_2(\Z))$。因为 $f$ 在 $\infty$ 处至少有一阶零点，而 $\Delta$ 在 $\infty$ 处恰有一阶零点，并且在上半平面中无零点，所以 $f/\Delta$ 仍在整个模曲线上全纯。其权重降低为 $k-12$，故
\[
\frac{f}{\Delta}\in M_{k-12}(\SL_2(\Z)).
\]
所以 $f=\Delta g$。

反过来，若 $g\in M_{k-12}(\SL_2(\Z))$，则 $\Delta g$ 显然是权 $k$ 尖点形式。因此
\[
S_k(\SL_2(\Z))=\Delta\,M_{k-12}(\SL_2(\Z))
\]
对每个 $k$ 都成立。

唯一性来自 $\Delta\neq 0$。故整个尖点理想恰为主理想 $(\Delta)$。
\end{solution}

\begin{problem}{模形式环由两个 Eisenstein 级数生成}
\hard
证明每个偶权全纯模形式都可写成 $E_4$ 与 $E_6$ 的多项式。
\end{problem}
\begin{solution}
对权 $k$ 作归纳。

当 $k=0,4,6,8,10$ 时，结论显然成立：这几个空间都是一维，生成元分别可取
\[
1,\ E_4,\ E_6,\ E_4^2,\ E_4E_6.
\]
而 $M_2(\SL_2(\Z))=0$。

设现在 $k\ge 12$ 且结论对更低偶权都成立。取任意 $f\in M_k(\SL_2(\Z))$。由于方程 $4a+6b=k$ 总有解，可取一个权 $k$ 的单项式
\[
P_k(E_4,E_6)=E_4^aE_6^b
\]
其常数项为 $1$。于是
\[
f-a_0(f)P_k(E_4,E_6)
\]
在 $\infty$ 处常数项为零，所以是尖点形式。由上一题，存在 $g\in M_{k-12}(\SL_2(\Z))$ 使
\[
f-a_0(f)P_k(E_4,E_6)=\Delta g.
\]
再用
\[
\Delta=\frac{E_4^3-E_6^2}{1728}
\]
以及归纳假设可知右边是 $E_4,E_6$ 的多项式，因此 $f$ 也是 $E_4,E_6$ 的多项式。

归纳完成，所以
\[
M_*(\SL_2(\Z))=\C[E_4,E_6].
\]
\end{solution}

\begin{problem}{两个生成元代数无关}
\hard
证明 $E_4$ 与 $E_6$ 代数无关。
\end{problem}
\begin{solution}
固定偶权 $k$，考虑所有满足 $4a+6b=k$ 的单项式 $E_4^aE_6^b$。在前一题中我们已经知道这些单项式张成整个 $M_k(\SL_2(\Z))$。

另一方面，这样的单项式个数恰好等于 $\dim M_k(\SL_2(\Z))$，这是前面单项式计数题已经证明过的。于是这些单项式不只是张成，而且构成一组基。

现在若存在非零多项式关系
\[
P(E_4,E_6)=0,
\]
把它按权分次分解后，可得到某个固定权里的非平凡线性关系，与“该权单项式构成基”矛盾。

因此不存在非零多项式关系，$E_4$ 与 $E_6$ 代数无关。
\end{solution}

\begin{problem}{素数水平中的亏格零与亏格一}
\hard
对素数 $p=5,7,11,13,17,19,23$，判定 $X_0(p)$ 的亏格是 $0$、$1$ 还是至少 $2$。
\end{problem}
\begin{solution}
对素数 $p>3$，有
\[
g(X_0(p))=\frac{p+1}{12}-\frac{\nu_2}{4}-\frac{\nu_3}{3}.
\]
逐个代入：
\begin{itemize}
  \item $p=5$：$\nu_2=2$、$\nu_3=0$，故 $g=\frac{6}{12}-\frac{1}{2}=0$。
  \item $p=7$：$\nu_2=0$、$\nu_3=2$，故 $g=\frac{8}{12}-\frac{2}{3}=0$。
  \item $p=11$：$\nu_2=0$、$\nu_3=0$，故 $g=\frac{12}{12}=1$。
  \item $p=13$：$\nu_2=2$、$\nu_3=2$，故 $g=\frac{14}{12}-\frac{1}{2}-\frac{2}{3}=0$。
  \item $p=17$：$\nu_2=2$、$\nu_3=0$，故 $g=\frac{18}{12}-\frac{1}{2}=1$。
  \item $p=19$：$\nu_2=0$、$\nu_3=2$，故 $g=\frac{20}{12}-\frac{2}{3}=1$。
  \item $p=23$：$\nu_2=0$、$\nu_3=0$，故 $g=\frac{24}{12}=2$。
\end{itemize}
因此：
\[
\text{亏格 }0:\ p=5,7,13;
\qquad
\text{亏格 }1:\ p=11,17,19;
\qquad
\text{至少 }2:\ p=23.
\]
\end{solution}

\begin{problem}{一级别十一的整条维数公式}
\hard
证明对所有偶数 $k\ge 4$，有
\[
\dim M_k(\Gamma_0(11))=k,
\qquad
\dim S_k(\Gamma_0(11))=k-2.
\]
\end{problem}
\begin{solution}
对 $\Gamma_0(11)$，几何数据是
\[
g=1,
\qquad \nu_2=0,
\qquad \nu_3=0,
\qquad \nu_\infty=2.
\]
代入一般维数公式，得到
\[
\dim M_k(\Gamma_0(11))=(k-1)(1-1)+0+0+\frac{k}{2}\cdot 2=k.
\]
对偶数 $k\ge 4$，尖点空间维数是
\[
\dim S_k(\Gamma_0(11))=(k-1)(1-1)+0+0+\left(\frac{k}{2}-1\right)\cdot 2=k-2.
\]
这就得到整条线性的维数公式。

顺便检查 $k=2$：此时
\[
\dim S_2(\Gamma_0(11))=g=1,
\]
也与这一组数据一致。
\end{solution}

\begin{problem}{从局部参数看维数公式中的整数部分}
\hard
解释为什么一般维数公式里会出现
\[
\left\lfloor\frac{k}{4}\right\rfloor\nu_2,
\qquad
\left\lfloor\frac{k}{3}\right\rfloor\nu_3,
\qquad
\frac{k}{2}\nu_\infty
\]
这些项。
\end{problem}
\begin{solution}
核心思想是：权 $k$ 模形式在特殊点附近并不是普通函数，而要乘上相应的自动因子。

在阶 $2$ 椭圆点附近，真正的局部参数是平方根型的。若把模形式改写成局部坐标中的普通函数，可允许的极点阶最多是 $k/4$，而极点阶必须是整数，所以贡献是
\[
\left\lfloor\frac{k}{4}\right\rfloor.
\]

在阶 $3$ 椭圆点附近，局部参数是立方根型的。同理，允许的极点阶最多是 $k/3$，因此贡献变成
\[
\left\lfloor\frac{k}{3}\right\rfloor.
\]

在尖点附近，局部参数是 $q_h=e^{2\pi i\tau/h}$。权 $k$ 形式的增长由 $(d\tau)^{k/2}$ 控制，于是每个尖点贡献 $k/2$ 个自由度；若要求是尖点形式，则还要额外消失一次，所以每个尖点只剩
\[
\frac{k}{2}-1.
\]
这些局部贡献与 Riemann--Roch 的全局项 $(k-1)(g-1)$ 合在一起，就得到整条维数公式。
\end{solution}

\begin{problem}{权二十四空间中的 Hecke 分解现象}
\hard
说明为什么 $S_{24}(\SL_2(\Z))$ 已经不再是自动的一维 Hecke 本征空间，并解释“寻找 Hecke 基”在这一权里意味着什么。
\end{problem}
\begin{solution}
由维数公式，
\[
\dim S_{24}(\SL_2(\Z))=2.
\]
例如可以取一组基为
\[
\Delta E_4^3,
\qquad
\Delta^2.
\]
因此单凭维数已经不能像一维情形那样唯一确定特征形式。

Hecke 算子 $T_n$ 仍然两两交换，并且对 Petersson 内积自伴，所以它们可以在这个二维空间中同时对角化。于是存在一组基
\[
f_1,\ f_2
\]
使得对每个 $n$ 都有
\[
T_n f_j=\lambda_{n,j}f_j.
\]
这组基就叫作 Hecke 基。

所以在权 $24$ 的情形中，问题不再是“空间是否存在非零尖点形式”，而是“如何在二维尖点空间中选出真正的本征方向”。这正是高维 Hecke 理论开始变得有内容的地方。
\end{solution}
